\documentclass[10pt]{article}
\usepackage[a4paper,margin=0.55in]{geometry}
\usepackage{amsmath,amssymb,mathtools}
\usepackage{booktabs}
\usepackage{xcolor}
\usepackage{hyperref}
\setlength{\jot}{5pt}
\numberwithin{equation}{section}
\allowdisplaybreaks
\newcommand{\red}[1]{{\color{red}#1}}

\title{Limiting cases ($\Lambda\to0$, $\vee\to\infty$) of the finite terms:\\ cancellation of the negative-argument logarithms, reality, and simplification}
\date{}
\begin{document}
\maketitle

\section{Where an $i$ can come from}
In the limiting-case file the possible sources of an imaginary part are:
\begin{enumerate}
\item[(C1)] an explicit $i$ in the coefficient ($ig^{4}f^{abd}$, $ig^{4}f^{c'd'a}$, \dots);
\item[(C2)] the Fourier phases $e^{-ik\cdot(w'-w)}$, $e^{-ik\cdot(y'-w)}$, $e^{-ik\cdot(z-x)p^{+}/k^{+}}$, \dots;
\item[(C3)] the momentum integral of Group I, Row II ($N_{2b}$);
\item[(C4)] logarithms whose argument can be negative. In the limiting case these occur \emph{only} in Group I, Row III.
\end{enumerate}
All Wilson lines and $f^{abc}$ are real, $\rho$ is hermitian, every transverse kernel is a real rational function, and $\mathcal A_{0}$ is real. Its radicand $(x-w)^{2}(x-z)^{2}-[(x-w)\cdot(x-z)]^{2}\ge0$ by Cauchy--Schwarz.

\textbf{Rule.} Groups I and III are interferences $\mathcal A_{\rm LO}^{\dagger}\mathcal A_{\rm NLO}$ and come with their c.c.; Group II is $|\mathcal A|^{2}$ and must be hermitian by itself. For a term $T=\int e^{-ik\cdot\Delta}\,c\,\mathcal K\,\mathcal O$ with real kernel $\mathcal K$ and classical $\rho$:
\begin{equation}
c\ \text{real:}\quad T+T^{\dagger}=2c\int\cos(k\cdot\Delta)\,\mathcal K\,\mathcal O,\qquad
c=i\alpha:\quad T+T^{\dagger}=2\alpha\int\sin(k\cdot\Delta)\,\mathcal K\,\mathcal O,
\label{eq:rule}
\end{equation}
because $i\,e^{-ik\cdot\Delta}+(i\,e^{-ik\cdot\Delta})^{*}=2\sin(k\cdot\Delta)$.

\section{Logarithms in the limiting-case file}
\begin{tabular}{p{0.2\textwidth}p{0.45\textwidth}p{0.27\textwidth}}
\toprule
Row & logarithms & argument\\\midrule
GI R2 ($N_{2b}$) & $\log\frac{p^{2}}{k^{2}}$, $\log\frac{k^{2}}{(k-p)^{2}}$ & positive\\
GI R3 & $\log\frac{(x-z)^{2}}{(x-z)^{2}-(x-w)^{2}}$, $\log\frac{(x-w)^{2}-(x-z)^{2}}{(x-w)^{2}}$, $\log\frac{-(x-z)^{2}}{(x-w)^{2}}$, $\log\frac{-(x-z)^{2}}{(x-w)^{2}-(x-z)^{2}}$ & \red{can be negative}\\
GI R3 ($N_{3b}$) & $\log\frac{(z-w)^{2}}{(x-w)^{2}}$ & positive\\
GI R4, GII R1, GIII R1, R2 & $\log\frac{(x-z)^{2}}{(x-w)^{2}}$, $\log\frac{(x'-z)^{2}}{(x'-w')^{2}}$, \dots & positive\\
GII R1 ($N_{2d,2e,2f}$) & $\log\frac{(x-z)^{2}}{(x-w)^{2}}\big/[(x-z)^{2}-(x-w)^{2}]$ & positive; the pole is removable (the log vanishes there)\\\bottomrule
\end{tabular}

\noindent So the only logarithms with negative arguments are in Group I, Row III. Sec.~3 treats them in full.

\section{Group I, Row III}
\subsection{Notation}
\begin{equation}
X\equiv x-w,\quad Z\equiv x-z,\quad W\equiv z-w=X-Z,\quad a\equiv x'-w',\quad n\equiv y-z,\quad u\equiv\frac{Z^{2}}{X^{2}}=\frac{(x-z)^{2}}{(x-w)^{2}}>0 .
\end{equation}
All four pieces $N_{1},N_{2},N_{3a},N_{3b}$ of Row III have the \emph{same} colour structure
\begin{equation}
\mathcal C_{\rm III}=\Big[U^{ab'}(x')-U^{ab'}(w')\Big]\rho^{b'}(x')\,\Big[U^{de}(x)\,U^{bc}(z)-U^{de}(x)\,U^{bc}(y)\Big]\rho^{e}(x)\,\rho^{c}(y),
\end{equation}
and the same overall factor. So the whole row can be written as one kernel:
\begin{equation}
\frac{d^{3}N_{\rm III}}{d^{3}k}=\frac{1}{(2\pi)^{3}}\,\frac{ig^{4}f^{abd}}{16\pi^{5}}\,\frac{1}{k^{+}}\int_{w,w'}e^{-ik\cdot(w'-w)}\int_{x,y,z,x'}\frac{a^{i}\,n^{j}}{a^{2}\,n^{2}}\;M^{ij}(x,z,w)\;\mathcal C_{\rm III}.
\label{eq:unit}
\end{equation}
Terms written with $\frac{ig^{4}}{8\pi^{5}}$ contribute with a factor $2$ to $M^{ij}$.

\subsection{Step 1: the branch of the logarithms}
The logs come from $p^{+}$ integrals with the energy denominator $p^{+}(x-z)^{2}-k^{+}(x-w)^{2}$ of the single-gluon state $\bar{\mathbb B}_{3}$. In $\bar{\mathbb B}_{3}^{\dagger}$ this denominator carries $-i0$. That is exactly the \emph{principal branch}, $\log(-A)=\log A+i\pi$ for $A>0$, which was checked earlier at finite $\vee$ against the residue sum to 12 digits. The limit $\vee\to\infty$ ($\kappa=k^{+}/\vee\to0$) does not change the branch. So no absolute values are needed.

\subsection{Step 2: every log on the principal branch}
With $\theta(\cdot)$ the step function:
\begin{align}
L_{1a}&\equiv\log\frac{Z^{2}}{Z^{2}-X^{2}}=\log\frac{u}{u-1}=\log\frac{u}{|1-u|}+i\pi\,\theta(1-u),\\
L_{1b}&\equiv\log\frac{X^{2}-Z^{2}}{X^{2}}=\log(1-u)=\log|1-u|+i\pi\,\theta(u-1),\\
L_{1c}&\equiv\log\frac{-Z^{2}}{X^{2}}=\log(-u)=\log u+i\pi,\\
L_{2}&\equiv\log\frac{-Z^{2}}{X^{2}-Z^{2}}=\log\frac{-u}{1-u}=\log\frac{u}{|1-u|}+i\pi\,\theta(1-u),\\
L_{3}&\equiv\log\frac{X^{2}-Z^{2}}{X^{2}}=\log(1-u)=\log|1-u|+i\pi\,\theta(u-1).
\end{align}
Adding the pairs that appear together, the $\log|1-u|$ and the step functions cancel:
\begin{equation}
\boxed{\;L_{1a}+L_{1b}=L_{2}+L_{3}=L_{1c}=\log u+i\pi=\log(-u)\qquad\text{for all }u>0.\;}
\label{eq:logsum}
\end{equation}
For example, for $u<1$: $L_{2}+L_{3}=\big[\log\frac{u}{1-u}+i\pi\big]+\log(1-u)=\log u+i\pi$. For $u>1$: $L_{2}+L_{3}=\log\frac{u}{u-1}+\big[\log(u-1)+i\pi\big]=\log u+i\pi$.

\subsection{Step 3: collecting the logarithmic terms}
$N_{1}$: the first two terms have the same kernel, so they give $(L_{1a}+L_{1b})$. The third term has $L_{1c}$:
\begin{equation}
M^{ij}_{N_{1}}=\frac{X^{2}+Z^{2}}{(Z^{2}-X^{2})^{2}}\,\delta^{ij}\,(L_{1a}+L_{1b})-\frac{2X^{i}Z^{j}}{X^{2}Z^{2}}\,L_{1c}.
\end{equation}
$N_{2}$ (all with $L_{2}$, overall sign $-$, factor 2) and the first six terms of $N_{3a}$ (all with $L_{3}$, factor 2) contain the \emph{same six tensors}:
\begin{center}
\begin{tabular}{lll}
\toprule
tensor & $N_{2}$ & $N_{3a}$\\\midrule
$\delta^{ij}\frac{(W\cdot X)Z^{2}}{W^{2}(X^{2}-Z^{2})^{2}}$ & $-2L_{2}$ & $-2L_{3}$\\
$\delta^{ij}\frac{(W\cdot Z)X^{2}}{W^{2}(X^{2}-Z^{2})^{2}}$ & $+2L_{2}$ & $+2L_{3}$\\
$\frac{W^{i}X^{j}}{W^{2}(X^{2}-Z^{2})}$ & $-2L_{2}$ & $-2L_{3}$\\
$\frac{X^{i}W^{j}Z^{2}}{W^{2}(X^{2}-Z^{2})X^{2}}$ & $+2L_{2}$ & $+2L_{3}$\\
$\frac{W^{i}Z^{j}X^{2}}{W^{2}(X^{2}-Z^{2})Z^{2}}$ & $+2L_{2}$ & $+2L_{3}$\\
$\frac{Z^{i}W^{j}}{W^{2}(X^{2}-Z^{2})}$ & $-2L_{2}$ & $-2L_{3}$\\\bottomrule
\end{tabular}
\end{center}
Hence they combine into $(L_{2}+L_{3})$. Using \eqref{eq:logsum}, all logarithmic terms of $N_{1}+N_{2}+N_{3a}$ share the single factor $\log(-u)$:
\begin{align}
M^{ij}_{\log}&=\log(-u)\Bigg\{\frac{X^{2}+Z^{2}}{(X^{2}-Z^{2})^{2}}\,\delta^{ij}-\frac{2X^{i}Z^{j}}{X^{2}Z^{2}}
+\frac{2}{W^{2}(X^{2}-Z^{2})}\Big[\frac{(W\cdot Z)X^{2}-(W\cdot X)Z^{2}}{X^{2}-Z^{2}}\,\delta^{ij}\notag\\
&\qquad-W^{i}X^{j}+\frac{Z^{2}}{X^{2}}X^{i}W^{j}+\frac{X^{2}}{Z^{2}}W^{i}Z^{j}-Z^{i}W^{j}\Big]\Bigg\}.
\end{align}

\subsection{Step 4: the bracket collapses to the Row IV kernel}
\textbf{$\delta^{ij}$ part.} With $W=X-Z$: $W\cdot Z=X\cdot Z-Z^{2}$, $W\cdot X=X^{2}-X\cdot Z$, $W^{2}=X^{2}+Z^{2}-2X\cdot Z$, so
\[
(W\cdot Z)X^{2}-(W\cdot X)Z^{2}=(X\cdot Z)(X^{2}+Z^{2})-2X^{2}Z^{2},
\]
\[
\frac{(X^{2}+Z^{2})W^{2}+2(X\cdot Z)(X^{2}+Z^{2})-4X^{2}Z^{2}}{W^{2}(X^{2}-Z^{2})^{2}}=\frac{(X^{2}+Z^{2})^{2}-4X^{2}Z^{2}}{W^{2}(X^{2}-Z^{2})^{2}}=\frac{1}{W^{2}} .
\]
\textbf{Non-$\delta$ part.} One needs
$\frac{1}{X^{2}-Z^{2}}\big[-W^{i}X^{j}+\frac{Z^{2}}{X^{2}}X^{i}W^{j}+\frac{X^{2}}{Z^{2}}W^{i}Z^{j}-Z^{i}W^{j}\big]=\frac{W^{i}Z^{j}}{Z^{2}}-\frac{X^{i}W^{j}}{X^{2}}$.
Multiply by $X^{2}-Z^{2}$; the difference of the two sides is $-W^{i}(X-Z)^{j}+(X-Z)^{i}W^{j}=-W^{i}W^{j}+W^{i}W^{j}=0$. Therefore
\begin{equation}
\boxed{\;M^{ij}_{\log}=2\,\log(-u)\;\mathcal K^{ij}_{\rm IV},\qquad
\mathcal K^{ij}_{\rm IV}=\frac{\delta^{ij}}{2(z-w)^{2}}+\frac{(z-w)^{i}(x-z)^{j}}{(z-w)^{2}(x-z)^{2}}-\frac{(x-w)^{i}(z-w)^{j}}{(z-w)^{2}(x-w)^{2}}-\frac{(x-w)^{i}(x-z)^{j}}{(x-w)^{2}(x-z)^{2}},\;}
\end{equation}
which is exactly the kernel of Group I, Row IV. All the $\frac{1}{(x-z)^{2}-(x-w)^{2}}$ poles have disappeared. \textbf{Numerical check:} for random $x,z,w$ (both $u<1$ and $u>1$), $\max|M^{ij}_{\log}-2(\log u+i\pi)\mathcal K^{ij}_{\rm IV}|\le7\times10^{-13}$, computed from the file's terms on the principal branch.

\subsection{Step 5: the $\mathcal A_{0}$ terms}
Since $Z^{2}-X\cdot Z=-Z\cdot W$ and $|X\times Z|=|W\times Z|$, the two arctangents add up to the angle of the triangle $(x,z,w)$ at $w$:
\begin{equation}
\boxed{\;\mathcal A_{0}=\frac{\theta_{w}}{\big|(x-w)\times(z-w)\big|},\qquad \theta_{w}=\arccos\frac{(x-w)\cdot(z-w)}{|x-w|\,|z-w|}\in[0,\pi],\;}
\end{equation}
which was checked numerically to $10^{-12}$. The four $\mathcal A_{0}$ terms ($N_{3a}$: 3, $N_{3b}$: 1) combine, using the projections $X_{\perp}=X-\frac{(X\cdot Z)}{Z^{2}}Z$ and $Z_{\perp}=Z-\frac{(X\cdot Z)}{X^{2}}X$, into
\begin{equation}
a^{i}n^{j}M^{ij}_{\mathcal A}=\frac{2\,\mathcal A_{0}}{(z-w)^{2}}\Big[(a\cdot W)\,(n\cdot X_{\perp})+(a\cdot Z_{\perp})\,(n\cdot W)\Big].
\end{equation}

\subsection{Step 6: the remaining $\log\frac{(z-w)^{2}}{(x-w)^{2}}$ terms ($N_{3b}$)}
\begin{equation}
a^{i}n^{j}M^{ij}_{B}=\frac{1}{(z-w)^{2}}\log\frac{(z-w)^{2}}{(x-w)^{2}}\Big[-\frac{(a\cdot W)(n\cdot Z)}{Z^{2}}+\frac{(a\cdot X)(n\cdot W)}{X^{2}}-a\cdot n\Big].
\end{equation}

\subsection{Step 7: Row III, simplified}
\begin{align}
\frac{d^{3}N_{\rm III}}{d^{3}k}&=\frac{1}{(2\pi)^{3}}\,\frac{ig^{4}f^{abd}}{16\pi^{5}}\,\frac{1}{k^{+}}\int_{w,w'}e^{-ik\cdot(w'-w)}\int_{x,y,z,x'}\frac{1}{(x'-w')^{2}(y-z)^{2}}\Bigg\{\notag\\&2\log\frac{(x-z)^{2}}{(x-w)^{2}}\,(x'-w')^{i}(y-z)^{j}\,\mathcal K^{ij}_{\rm IV}\notag\\
&+\frac{2\,\theta_{w}}{|(x-w)\times(z-w)|\,(z-w)^{2}}\Big[(x'-w')\cdot(z-w)\;(y-z)\cdot X_{\perp}+(x'-w')\cdot Z_{\perp}\;(y-z)\cdot(z-w)\Big]\notag\\
&+\frac{1}{(z-w)^{2}}\log\frac{(z-w)^{2}}{(x-w)^{2}}\notag\\&\quad\times\Big[-\frac{(x'-w')\cdot(z-w)\;(y-z)\cdot(x-z)}{(x-z)^{2}}+\frac{(x'-w')\cdot(x-w)\;(y-z)\cdot(z-w)}{(x-w)^{2}}-(x'-w')\cdot(y-z)\Big]\Bigg\}\mathcal C_{\rm III}\notag\\
&+\frac{d^{3}N_{i\pi}}{d^{3}k},
\label{eq:R3final}
\end{align}
with the imaginary part of $\log(-u)$, multiplied by the overall $i$, giving the \emph{real} term
\begin{equation}
\frac{d^{3}N_{i\pi}}{d^{3}k}=-\frac{1}{(2\pi)^{3}}\,\frac{g^{4}f^{abd}}{8\pi^{4}}\,\frac{1}{k^{+}}\int_{w,w'}e^{-ik\cdot(w'-w)}\int_{x,y,z,x'}\frac{(x'-w')^{i}(y-z)^{j}}{(x'-w')^{2}(y-z)^{2}}\,\mathcal K^{ij}_{\rm IV}\;\mathcal C_{\rm III}.
\end{equation}
(At finite $\vee$ this carried $\theta(u-\kappa)$; for $\vee\to\infty$, $\theta(u)=1$.)

\subsection{Step 8: adding the c.c.}
The first three lines of \eqref{eq:R3final} have coefficient $i\times$real, so they are class $i\alpha$ with $\alpha=\frac{g^{4}f^{abd}}{16\pi^{5}k^{+}}$. $N_{i\pi}$ has a real coefficient. By \eqref{eq:rule}, for classical $\rho$:
\begin{align}
\frac{d^{3}N_{\rm III}}{d^{3}k}+{\rm c.c.}&=\frac{1}{(2\pi)^{3}}\,\frac{g^{4}f^{abd}}{8\pi^{5}}\,\frac{1}{k^{+}}\int\sin\big(k\cdot(w'-w)\big)\,\Big\{\text{first three lines of \eqref{eq:R3final}}\Big\}\,\mathcal C_{\rm III}\notag\\
&\quad-\frac{1}{(2\pi)^{3}}\,\frac{g^{4}f^{abd}}{4\pi^{4}}\,\frac{1}{k^{+}}\int\cos\big(k\cdot(w'-w)\big)\,\frac{(x'-w')^{i}(y-z)^{j}}{(x'-w')^{2}(y-z)^{2}}\,\mathcal K^{ij}_{\rm IV}\,\mathcal C_{\rm III},
\end{align}
which is manifestly real. No $\log$ of a negative number and no explicit $i$ is left.

\subsection{Step 9: singularities}
\begin{itemize}
\item \textbf{$(x-z)^{2}=(x-w)^{2}$.} All such poles cancelled in Step 4.
\item \textbf{$z\to y$, $x\to z$, $x\to w$.} All are integrable: at most $1/|\cdot|$ or $\log|\cdot|/|\cdot|$.
\item \textbf{$z\to w$, in $N_{i\pi}$.} $\mathcal K_{\rm IV}\to\frac{\delta^{ij}}{2(z-w)^{2}}$, which is log divergent. At $z=w$ the colour structure becomes $f^{abd}X^{a}(w')\cdot Y^{b}(w)\,D^{d}$, with $X^{a}(w')=\int_{x'}\frac{(x'-w')}{(x'-w')^{2}}[U(x')-U(w')]^{ab'}\rho^{b'}(x')$, $Y^{b}(w)=\int_{y}\frac{(y-w)}{(y-w)^{2}}[U(w)-U(y)]^{bc}\rho^{c}(y)$ and $D^{d}=\int_{x}U^{de}(x)\rho^{e}(x)$. This is antisymmetric under $w\leftrightarrow w'$, and $N_{i\pi}$ has a real coefficient, so by \eqref{eq:rule} its divergent part is multiplied by $\cos(k\cdot(w'-w))$, which is symmetric, and vanishes. It is a \textbf{Coulomb phase}, and it cancels exactly against the c.c. The rest of $N_{i\pi}$ is finite.
\item \textbf{$z\to w$, in the $\log u$ and $\mathcal A_{0}$ terms.} Both are finite: $\log u=O(|z-w|)$ is odd in $z-w$, and $\mathcal A_{0}$ times its bracket is $O(1/|z-w|)$. Numerically, $r^{2}\times$(angle average) goes to $0$ as $r=|z-w|\to0$.
\item \red{\textbf{$z\to w$, in the last term of $N_{3b}$}: $-\frac{(x'-w')\cdot(y-z)}{(x'-w')^{2}(y-z)^{2}}\frac{1}{(z-w)^{2}}\log\frac{(z-w)^{2}}{(x-w)^{2}}$.} \red{This behaves like $\frac{\log r^{2}}{r^{2}}$, a \emph{double-log} divergence of $\int d^{2}z$. Numerically, $r^{2}\times$(angle average) $=2.98,\,5.10,\,7.22,\,9.34$ for $r=10^{-1},\dots,10^{-4}$: it grows by $\approx2\log10$ per decade. It is class $i\alpha$, so with the c.c.\ it multiplies $\sin(k\cdot(w'-w))$, which is antisymmetric, times the antisymmetric colour structure above. The product is symmetric, so it does \emph{not} cancel against the c.c. Nothing else in Row III has this $z\to w$ behaviour.}
\end{itemize}

\red{\textbf{Consequence.} With this one exception, Row III in the limit $\Lambda\to0$, $\vee\to\infty$ is finite and real. The $\frac{1}{(z-w)^{2}}\log\frac{(z-w)^{2}}{(x-w)^{2}}$ term of $N_{3b}$ cannot be taken to $\Lambda\to0$ on its own. It is the remainder of a $p^{+}$ integral whose $\log(k^{+}/\Lambda)\frac{1}{(z-w)^{2}}$ partner went into the JIMWLK file, and the two must be combined, with the same colour structure, before $\Lambda\to0$. Equivalently, the $z\to w$ region must be cut off at the scale set by $\Lambda$, e.g.\ $(z-w)^{2}\gtrsim\frac{\Lambda}{k^{+}}(x-w)^{2}$. The origin of this term needs to be rechecked.}

\section{The other rows}
\begin{tabular}{p{0.2\textwidth}p{0.12\textwidth}p{0.6\textwidth}}
\toprule
Row & class & how it is real\\\midrule
GI R2, $N_{2b}$ & $i\alpha$ (implicit) & $\Phi^{i}=\int d^{2}q\,d^{2}p\,e^{iq\cdot(w-y)+ip\cdot(y-x)}\frac{R^{i}(q,p)}{q^{2}p^{2}}$ is $i\times$real ($R^{i}$ is odd under $(q,p)\to(-q,-p)$). With c.c.: $\sin$. \red{Rename the dummy $\int d^{2}k\to\int d^{2}q$: it clashes with the external $k$.}\\
GI R4 & $i\alpha$ & positive log; with c.c.: $\frac{g^{4}}{4\pi^{5}}\sin(k\cdot(w'-w))$\\
GI R5, first part & $i\alpha$ & phase $e^{-ik\cdot(w'-Y_{p})}$ with real $Y_{p}=(1-\frac{p^{+}}{k^{+}})z+\frac{p^{+}}{k^{+}}x$; with c.c.: $\sin$\\
GI R5, $N_{2a}$, $N_{2b\,\rm rest}$ & real & with c.c.: $\cos$\\
GII R1, first part & --- & $N_{1b}=N_{1a}^{\dagger}$, so $N_{1a}+N_{1b}=-\frac{g^{4}}{8\pi^{5}k^{+}}\int\sin(k\cdot(w'-w))(\dots)$\\
GII R1, $N_{2a}$--$N_{2f}$ & real & each kernel is even under $(x,w)\leftrightarrow(x',w')$; hermitian, $\cos$\\
GII R4 & real & symmetric under $(x,y)\leftrightarrow(x',y')$; $\cos(k\cdot(y'-y)/\xi)$\\
GIII R1, R2, R3 (first), R5, R6 & $i\alpha$ & positive logs and denominators; with c.c.: $\sin$\\
GIII R3 (second) & real & with c.c.: $\cos$\\\bottomrule
\end{tabular}

\section{Summary}
\begin{enumerate}
\item \textbf{Negative-argument logs (only GI R3).} On the principal branch, which is the physical $-i0$ of $\bar{\mathbb B}_{3}^{\dagger}$, they combine \emph{exactly} into one factor $\log(-u)=\log\frac{(x-z)^{2}}{(x-w)^{2}}+i\pi$, multiplying $2\mathcal K_{\rm IV}$. All the $\log|1-u|$ pieces and step functions cancel, and so do all the $\frac{1}{(x-z)^{2}-(x-w)^{2}}$ poles.
\item \textbf{The $i\pi$.} Times the overall $i$ it becomes the \emph{real} term $N_{i\pi}$. Its only divergence ($z\to w$) is a Coulomb phase, which cancels against the c.c.
\item \textbf{Explicit $i$'s and phases.} They give $\sin$ (class $i\alpha$) or $\cos$ (real class) with the c.c.\ in Groups I and III. Group II is hermitian by itself. The final cross section is real.
\item \textbf{Simplifications.} Row III is eq.~\eqref{eq:R3final}, with $\mathcal A_{0}=\theta_{w}/|(x-w)\times(z-w)|$.
\item \red{\textbf{Open issue.} The last term of $N_{3b}$ has a non-integrable $\frac{\log(z-w)^{2}}{(z-w)^{2}}$ at $z\to w$ that survives the c.c. It must be recombined with its $\log(k^{+}/\Lambda)$ partner before $\Lambda\to0$.}
\end{enumerate}
\end{document}
